% Le titre IV (salaire) est à venir : on aligne la numérotation des sections
% sur celle des titres.
\setcounter{section}{4}
\section{Titre V --- Sources et hiérarchie des normes}

\article{Niveaux d'impérativité}\label{art:niveaux}
I. --- Chaque disposition du présent code est absolue, relative ou
dispositive. Son niveau est indiqué en tête de l'article. Un article sans
niveau ne prescrit rien : il définit une notion employée ailleurs.

II. --- Pour chaque disposition chiffrée, le code fixe le sens favorable au
salarié : une valeur plus grande ou une valeur plus petite. Le tableau
ci-dessous les récapitule.
\lean{Disposition}
\lean{Sens}

\article[absolue]{Dispositions absolues}\label{art:absolues}
Il ne peut être dérogé à une disposition absolue, ni par accord d'entreprise,
ni par contrat, même en faveur du salarié. Toute stipulation contraire est
réputée non écrite.
\lean{applicable}
\theoreme{applicable_absolue}

\article[absolue]{Dispositions relatives}\label{art:relatives}
I. --- Un accord d'entreprise ou le contrat ne peut déroger à une disposition
relative qu'en faveur du salarié.

II. --- Lorsque la loi, un accord d'entreprise et le contrat fixent chacun une
valeur, la valeur la plus favorable au salarié s'applique. Une stipulation
moins favorable que la loi est sans effet.

III. --- Le préavis dû par l'employeur ne peut qu'être allongé ; le préavis dû
par le salarié ne peut qu'être raccourci.
\lean{applicable}
\theoreme{applicable_relative_plancher}
\theoreme{applicable_relative_faveur}
\theoreme{applicable_relative_clause_defavorable}
\theoreme{preavisEmployeur_ge_legal}
\theoreme{preavisSalarie_le_legal}

\article[absolue]{Dispositions dispositives et primauté de l'accord d'entreprise}\label{art:dispositives}
I. --- Une disposition dispositive s'applique à défaut de stipulation
contraire.

II. --- L'accord d'entreprise l'emporte sur le contrat, et le contrat sur la
loi.

III. --- Lorsqu'une disposition dispositive comporte une limite absolue, une
stipulation qui la dépasse est réduite à cette limite.
\lean{applicable}
\theoreme{applicable_dispositive_accord}
\theoreme{applicable_dispositive_limite}
\theoreme{dureeEssai_le_max}

\begin{center}
\footnotesize
\begin{tabular}{llccc}
  \multicolumn{5}{c}{\textit{Dispositions chiffrées du code}} \\[0.3em]
  \hline
  Disposition & Niveau & Valeur légale & Sens favorable & Limite \\
  \hline
  Préavis dû par l'employeur & relative & 7 à 90 jours (art.~\ref{art:resiliation}) & plus long & --- \\
  Préavis dû par le salarié & relative & 7 à 90 jours (art.~\ref{art:resiliation}) & plus court & --- \\
  Durée de l'essai & dispositive & 30 jours & plus courte & 90 jours \\
  Indemnité de rupture & relative & formule (art.~\ref{art:indemnite}) & plus élevée & --- \\
  Durée maximale & absolue & 48 h par semaine en moyenne & --- & --- \\
  Journée du travailleur de nuit & absolue & 8 h & --- & --- \\
  Période de référence & dispositive & 17 semaines & plus courte & 26 semaines \\
  Pause & relative & 20 minutes & plus longue & --- \\
  Majoration des heures supplémentaires & dispositive & 10 \% & plus élevée & 0 \% \\
  Congés payés & relative & 20 jours ouvrés & plus nombreux & --- \\
  \hline
\end{tabular}
\end{center}
